\section{模型假设}
为简化问题，本文做出以下假设：
\begin{itemize}
	\item \textbf{假设1（单作满块）}：同一地块同一季次只种植一种作物，且作物种植面积填满整块地，不进行部分面积种植。
	\item \textbf{假设2（销售基准）}：各作物各季次的预期销售量以2023年实际产量为基准；2023年未生产或未销售的作物季次视为无销售基准，不纳入规划。
	\item \textbf{假设3（价格成本稳定）}：问题1中，各作物销售单价与种植成本取2023年附件给出的数值（价格取区间中点），且不随年份变化；亩产量取2023年值。
	\item \textbf{假设4（不确定性分布）}：问题2、3中，销售量、亩产量、价格的不确定性用截断标准正态随机变量刻画（各作物冲击幅度均落在题目给定波动区间内），情景间相互独立（问题2）或服从给定相关结构（问题3）；成本按每年5\%确定性增长。
	\item \textbf{假设5（忽略市场反馈）}：按题目约定，销售按预期销售量上限进行，不考虑价格随产量的内生调节，仅按题目给出的降价、折价规则处理超产部分。
	\item \textbf{假设6（地力轮作）}：同一地块连续两年同一季次不种同一种作物（避免重茬）；每个滚动三年窗口（2023--25、2024--26、……、2028--30）内至少种植一次豆类作物。
\end{itemize}

\section{符号说明}
本文使用的主要符号及其含义如表\ref{tab:symbols}所示。

\begin{table}[htbp]
	\caption{符号说明}
	\label{tab:symbols}
	\centering
	\begin{tabular}{cl}
		\toprule
		符号 & 说明 \\
		\midrule
		$\mathcal{P}$ & 地块集合，共54块；$\mathcal{P}_D,\mathcal{P}_F$为水浇地与智慧大棚 \\
		$\mathcal{T}$ & 年份集合 $\{2024,\dots,2030\}$ \\
		$\mathcal{S}$ & 季次集合 $\{\text{第一季},\text{第二季}\}$ \\
		$\mathcal{C}$ & 作物集合，共41种；$\mathcal{C}^b$为豆类作物集合 \\
		$\mathcal{C}_{ps}$ & 地块$p$季次$s$可种植的作物集合 \\
		$A_p$ & 地块$p$的面积（亩） \\
		$y_{c,p,s}$ & 作物$c$在地块$p$季次$s$的亩产量（斤/亩） \\
		$r_{c,s}$ & 作物$c$季次$s$的销售单价（元/斤，取2023区间中点） \\
		$c_{c,p,s}$ & 作物$c$在地块$p$季次$s$的种植成本（元/亩） \\
		$S_{c,s}$ & 作物$c$季次$s$的预期销售量（斤，=2023年产量） \\
		$x_{p,t,s,c}$ & 0-1决策变量：地块$p$第$t$年第$s$季是否种植作物$c$ \\
		$u_{\omega,c,s,t}$ & 情景$\omega$下第$t$年作物$c$季次$s$按原价售出的数量（斤） \\
		$v_{c,s,t}$ & 问题1情形2中第$t$年按50\%降价售出的数量（斤） \\
		$\Pi_\omega$ & 情景$\omega$下的总净收益（元） \\
		$\lambda$ & 风险偏好系数；$\alpha$为CVaR置信水平（取0.90） \\
		$\Omega$ & 随机情景集合，$|\Omega|=N$ \\
		\bottomrule
	\end{tabular}
\end{table}
